\section{Sleeping Through the Junta}

\begin{quotex}
Take heed, watch and pray, for you do not know when the time will come.
\flright{\textsc{James 13:33}}
\end{quotex}


\paragraph{The Movie Z}


The 1969 film ``Z'', about the military junta that took over Greece in the 60s, made quite an impact when it was released. A junta of colonels established a ``dictatorship'' in 1967, a month before a scheduled election, because they opposed the leftist policies of the likely election winner. The director of the film Z, Costas-Gavras, was adamantly opposed to the junta.

The anti-communist junta promoted conservative social policies. Anyone who opposed them were canceled, including journalists, artists, musicians, writers, philosophers, and even modern mathematics. When that failed, they resorted to the imprisonment, torture, and exile of political opponents. Crimes against such figures were covered up by government agencies.

\textbf{More astoundingly, they even banned the letter ``Z''!}


The Greek junta could not be missed as they were publicly visible. That tactic is not used much in our day. Instead, it can only be recognized by its effects. Neither the visible leaders, nor the purported political stance, matter. When journalists, writers, intellectuals, philosophers, and so on, are ``canceled'' then that is an indication of the junta in charge, not matter how or by whom. When political opponents are imprisoned and allies are not prosecuted, then the junta is in charge.

If could happen someday, so take heed, watch, and pray. When the junta takes control, let it not find you asleep.

\paragraph{Good and Evil at Play}


Religious people, of a certain sort, often claim that the world is ``good''. They claim the story of Genesis as the source, but that valuation applies only to the Edenic world. The expulsion of Adam and Eve from the Garden was supposed to be a punishment, not a vacation in the good world. Neither is the world wholly evil; rather it is a mixture. Hence, debates about conflicts in terms of good and evil are always fruitless.

The Law of Correspondence provides a better tool for analysis. The levels of being correspond to teach other in a way appropriate to that level. For example, here is a schema for the fundamental levels:

\begin{itemize}


\item \textbf{Spirit Level}: The Good


\item \textbf{Soul Level}: Good and Evil


\item \textbf{Material Level}: Friend and Enemy

\end{itemize}


Since the physical earth is in darkness, there is no battle between good and evil. Matter is neither good nor bad. The lion is not evil and the gazelle is not good. The distinction makes no sense. The one that does is the distinction between friend and enemy. The lion is the enemy to the gazelle, but there is no battle between good and evil.

Each combatant in a conflict has the right to self-defence, but that is not a battle between good and evil, except, perhaps, in a relative sense. But for spectators, is it really clear who is the friend and who is the enemy? Or is it better to stand aside?

It is different at the soul level, but how do you identify good and evil? There is talk in the air, in certain circles, about \emph{The Decline of the West}, \emph{The Crisis of the Modern World}, and
\emph{The Revolt Against the Modern World}. Herman Hesse, in \emph{A Glimpse into Chaos}, also argued that Europe is in a state of terminal decline.

Those works were written a century ago, so, if true, what is worth defending. Yet the times they are a -changing. Now the West is the leader of the free world, the protector of democracy, and on the right side of history. Anyone opposed to the agenda of the West is considered evil, and even canceled if necessary. On the other hand:

\begin{quotex}
The fulfilment of the ``American Dream'' is for many societies, especially traditional ones, an eschatological nightmare.
\flright{\textit{Catholic Herald}\footnote{\url{https://catholicherald.co.uk/putins-take-on-the-russian-idea/}}, 6 February 2022}
\end{quotex}


A decision must be made: Are the values of the West a sign of progress or an eschatological nightmare?

\url{https://www.youtube.com/watch?v=rHB28qfsV2w}

\flrightit{Posted on 2022-04-02 by Cologero}

\begin{center}* * *\end{center}

\begin{footnotesize}
\begin{sffamily}
\texttt{Santiago on 2022-04-04 at 19:27 said:}

So I recently went and saw Z.

Perhaps most conspicuous was the absence of any genuine political ideology throughout--towards the beginning of the film, the military junta describes their opposition as ``mildew'' which must be thoroughly eradicated through an impersonal process of ``heat'', ``cleansing'', etc. Such a materialistic, `scientistic' approach is hardly what you would expect from `far right' monarchists with self-professed Christian ideals.

The foremost general, Papadapoulos, would almost immediately--after the events of the film--opt for a reformed parliamentary system, perhaps to take the heat off of himself and his friends, which again is something you would not expect from a radical right-wing ideologue. His later move to formally abolish the monarchy all but confirms this.

The movie features some AntiF (radical left wing group) types of characters, who are in bed with the media, various workers institutions, even the higher ups of law enforcement, and funded by the political establishment. I thought it was revealing that 2 of this groups primary agents are both secretly Communists, and one is even a convicted child abuser, of the worst kind. Again, where is the motivating ideology behind the military junta if they are so readily willing to employ such people for their purposes?

The only distinction which is clearly apparent throughout, is that between ``friend'' and ``enemy.'' All ideology becomes nothing more than a vocabulary with which to dog-whistle to your ``friends'' that you are benevolent to the Regime.

American director William Friedkin said about the film: `` it was filmed like a documentary... Like the camera didn't know what was gonna happen next. And that is an induced technique.'' It's ironic now, because everything in that movie is, while tragic, utterly predictable. It's almost as if we're still banning letters of the alphabet today.

\hfill

\texttt{Greg on 2022-04-05 at 13:19 said:}

To answer your concluding question. The values of the west are indeed an eschatological nightmare. We are ruled by a high tech junta which will cancel you if you question their depraved values.

\hfill

\texttt{Tiffany on 2022-04-05 at 23:37 said:}

As an American, I would truly not be surprised if an alphabet letter were canceled this week. I can hear it now: Perhaps the letter `P' is ``triggering'' to some for its phallic nature, which promotes binary gender ``concepts'' and, I mean, its existence allows the word Patriarchy to exist(!).

I echo Greg's comment.

\end{sffamily}
\end{footnotesize}

